\begin{table}
\centering
\caption{\label{tbl-paper-conversion-complements}Predetermined Complements to Economic Conversion}
\centering
\fontsize{9}{11}\selectfont
\begin{threeparttable}
\begin{tabular}[t]{>{\raggedright\arraybackslash}p{4.0cm}>{\centering\arraybackslash}p{2.8cm}>{\centering\arraybackslash}p{2.8cm}>{\centering\arraybackslash}p{2.8cm}}
\toprule
Complement & All-Child EMI & Private EMI & Linguistic Distance\\
\midrule
\addlinespace[0.3em]
\multicolumn{4}{l}{\cellcolor{white}{\textbf{Panel A. Predetermined Census-2001 complements:}}}\\
\hspace{1em}Baseline human capital & 0.015** (0.004) & 0.015 (0.008) & \\
\hspace{1em}Urbanization & 0.018* (0.006) & 0.020 (0.010) & \\
\hspace{1em}ST concentration & 0.000 (0.005) & 0.004 (0.008) & \\
\addlinespace[0.3em]
\multicolumn{4}{l}{\cellcolor{white}{\textbf{Panel B. Predetermined IT environment:}}}\\
\hspace{1em}Baseline IT employment share & 0.001 (0.004) &  & 0.008 (0.010)\\
\bottomrule
\end{tabular}
\begin{tablenotes}[para]
\item \textit{Note: } 
\item Entries are interaction coefficients with state-clustered standard errors in parentheses; stars use Holm-adjusted p-values within each predeclared panel family. Panel A reports the change in the 2004-05 to 2022-23 schooling-consumption association for a one-standard-deviation increase in each predetermined Census-2001 complement, per 10 percentage points of schooling exposure. Panel B reports the corresponding 2005 IT-employment-environment interaction for all-child EMI per 10 percentage points and for linguistic distance per one Shastry degree. * p < 0.05, ** p < 0.01, *** p < 0.001 These are descriptive heterogeneity estimates and do not make the schooling exposure or linguistic-distance interaction causal.
\end{tablenotes}
\end{threeparttable}
\end{table}
